% command to make sure colorboxes don't create whitespace between code lines
\fboxsep=.6pt

% default settings for listings
\lstset{style=python}

\title[Daten Vorhersagen]{Daten Vorhersagen}

\subtitle{Markov-Ketten}

\author[Wendl, Cyril]{Cyril Wendl}
%\author[Wendl, Cyril]{Cyril Wendl\inst{2} \and Naoki Peter \inst{1}}

\institute[KZN] % (optional)
{
    % \inst{1}%
    % Fachschaft Informatik\\
    % Kantonsschule Zürich-Nord
    % \and
    % \inst{2}%
    Fachschaft Informatik\\
    Kantonsschule im Lee
}

\date[\today] % (optional)

\logo{\includegraphics[width=2cm]{\thelogo}}

\titlegraphic{\vspace{-.7cm}\includegraphics[width=\textwidth]{"GrundlagenInfo/06_Datenbanken/Figures/title-emojis"}}

\frame{\titlepage}

\section{Markov}

\begin{frame}{Themenübersicht}
    \begin{enumerate}
        \item[\faCheckSquare{}] Daten verstehen (analysieren und visualisieren)
        \item[\faCheckSquare{}] Zahlen vorhersagen (lineare Regression)
        \item[\faSquare] \textbf{Zustände vorhersagen (Markov-Ketten)}
    \end{enumerate}
\end{frame}

\begin{frame}{Anwendungen von Markov-Ketten}{Google PageRank-Algorithmus}
    \begin{itemize}[<+->]
        \item Das Surfen im Web kann als Markov-Kette modelliert werden.
        \item Jeder Kreis steht für eine Webseite (Zustand).
        \item Die Pfeile zeigen die Übergangswahrscheinlichkeiten zwischen den Seiten.
        \item Die Wahrscheinlichkeit, auf eine andere Seite zu wechseln, hängt nur von der aktuellen Seite ab – genau wie bei einer Markov-Kette!
    \end{itemize}

    \begin{figure}[H]
        \begin{tikzpicture}[scale=0.7, transform shape, ->, >=stealth', node distance=4cm, thick,
                website/.style={rectangle, draw, minimum width=2.5cm, minimum height=1cm, align=center}]

            \node[website] (A) {news.com};
            \node[website] (B) [right of=A] {blog.org};
            \node[website] (C) [below of=A] {shop.net};
            \node[website] (D) [below of=B] {wiki.edu};

            \path[->]
            (A) edge[bend left] node[above] {0.5} (B)
            (A) edge[bend right=15] node[left] {0.5} (C)
            (B) edge[bend left] node[below] {1.0} (A)
            (B) edge node[right] {1.0} (C)
            (C) edge node[below right] {1.0} (D)
            (D) edge[bend left] node[below] {0.6} (A)
            (D) edge[bend right=15] node[left] {0.4} (B);
        \end{tikzpicture}
    \end{figure}
\end{frame}

\begin{frame}{Anwendungen von Markov-Ketten}{Wortgenerierung mit ChatGPT}
    \begin{itemize}[<+->]
        \item Die Textgenerierung in ChatGPT kann als Markov-Kette modelliert werden.
        \item Jeder Kreis steht für ein Wort (Zustand).
        \item Die Pfeile zeigen die Übergangswahrscheinlichkeiten zum nächsten Wort.
        \item Die Wahrscheinlichkeit, welches Wort als nächstes kommt, hängt (im einfachsten Modell) nur vom aktuellen Wort ab – wie bei einer Markov-Kette!
    \end{itemize}

    \begin{figure}[H]
        \centering
        \begin{tikzpicture}[scale=0.7, transform shape, ->, >=stealth', node distance=4cm, thick,
                word/.style={rectangle, draw, minimum width=2.5cm, minimum height=1cm, align=center}]

            \node[word] (Hallo) {Hallo};
            \node[word] (wie) [right of=Hallo] {wie};
            \node[word] (geht) [below of=Hallo] {geht};
            \node[word] (dir) [below of=wie] {dir?};

            \path[->]
            (Hallo) edge[bend left] node[above] {0.8} (wie)
            (Hallo) edge[bend right=15] node[left] {0.2} (geht)
            (wie) edge node[right] {.6} (geht)
            (wie) edge node[right] {.4} (dir)
            (geht) edge node[below right] {1.0} (dir);
        \end{tikzpicture}
    \end{figure}
\end{frame}



\begin{frame}{Anwendungen von Markov-Ketten}
    \begin{columns}
        \column{0.5\textwidth}
        \begin{figure}[H]
            \centering
            \includegraphics[width=0.25\pagewidth]{"GrundlagenInfo/10_AusDatenLernen/Figures/intro_chatGPT"}
        \end{figure}
        \begin{itemize}[<+->]
            \item \textbf{Textgenerierung} in ChatGPT und anderen KI-Modellen
            \item \textbf{Wettervorhersage}
            \item \textbf{PageRank-Algorithmus} von Google
            \item \textbf{Biologische Systeme} (Gensequenzen)
        \end{itemize}

        \column{0.5\textwidth}
        \begin{itemize}[<+->]
            \item \textbf{Finanzmodelle} (Aktienkurse)
            \item \textbf{Verkehrssimulationen}
            \item \textbf{Warteschlangensysteme}
            \item \textbf{Spieltheorie} (Entscheidungsprozesse)
            \item \ldots
        \end{itemize}
    \end{columns}
\end{frame}

\begin{frame}[fragile]{Lernziele}
    \begin{itemize}[<+->]
        \item \faCheckSquare{} Ich kann \textbf{Anwendungen} von Markov-Ketten in der Informatik beschreiben.
        \item \faSquare{} Ich kann erklären, was eine \textbf{Markov-Kette} ist und wie die \textbf{Markov-Eigenschaft} funktioniert.
        \item \faSquare{} Ich kann aus einer realen Problemstellung eine \textbf{Übergangsmatrix} erstellen.
        \item \faSquare{} Ich kann mithilfe von Matrizen berechnen, wie sich eine Markov-Kette \textbf{über mehrere Schritte entwickelt}.
              % \item \faSquare{} Ich kann die stationäre Verteilung einer Markov-Kette berechnen und interpretieren.
        \item \faSquare{} Ich kann kurze Texte mit Markov-Ketten \textbf{generieren}.
        \item \faSquare{} Ich kann die \textbf{Grenzen} von Markov-Ketten reflektieren.
    \end{itemize}
\end{frame}

\begin{frame}{Was sind Markov-Ketten?}
    \begin{mydefinition}
    \begin{itemize}[<+->]
        \item Mathematisches Modell für \textbf{Abfolgen von Zuständen}
        \item \textbf{Markov-Eigenschaft}: Der nächste Zustand hängt \textbf{nur vom aktuellen Zustand} ab
    \end{itemize}
    \end{mydefinition}
\end{frame}

\begin{frame}
    \begin{myexercise}[Markov-Eigenschaft gegeben oder nicht?]
        \begin{enumerate}
            \item Das Wetter morgen kann gut vorhergesagt werden anhand des heutigen Wetters.
            \item Die Wahrscheinlichkeit, dass ein Schüler heute seine Hausaufgaben macht, hängt davon ab, ob er sie gestern und vorgestern gemacht hat.
            \item Die Position einer Spielfigur auf einem Spielbrett hängt nur vom aktuellen Feld und dem nächsten Würfelergebnis ab.
            \item Die Wahrscheinlichkeit, dass ein Computer abstürzt, hängt davon ab, was bereits alles auf dem Computer gemacht wurde.
            \item Die Wahrscheinlichkeit, dass eine Person Kleider einer bestimmten Marke kauft, hängt nur davon ab, ob das zuletzt gekaufte Kleidungsstück von dieser Marke war.
        \end{enumerate}
    \end{myexercise}
\end{frame}

\begin{frame}{Beispiel: Wettervorhersage}
    \begin{itemize}[<+->]
        \item Zwei mögliche Zustände: \emoji{sun} (sonnig) oder \emoji{umbrella} (regnerisch)
        \item Übergangswahrscheinlichkeiten:
              \begin{itemize}
                  \item Wenn heute \emoji{sun}: Morgen zu 80\% \emoji{sun}, zu 20\% \emoji{umbrella}
                  \item Wenn heute \emoji{umbrella}: Morgen zu 40\% \emoji{sun}, zu 60\% \emoji{umbrella}
              \end{itemize}
    \end{itemize}
    \only<5->{
        \begin{table}[H]
            \centering
            \begin{tabular}{c c|c|c}
                                                              &                  & \multicolumn{2}{c}{\textbf{Nach}}                    \\
                                                              &                  & \emoji{sun}                       & \emoji{umbrella} \\\midrule\midrule
                \multirow{2}{*}{\rotatebox{90}{\textbf{Von}}} & \emoji{sun}      & 0{,}8                             & 0{,}2            \\\cmidrule{2-4}
                                                              & \emoji{umbrella} & 0{,}4                             & 0{,}6            \\
            \end{tabular}
        \end{table}
    }
    \only<6->{
        \begin{myexercise}[Social-Media-Apps]\small
            Jedes Mal, wenn Sie Instagram benutzen, benutzen Sie als nächste App wieder Instagram mit 80\% Wahrscheinlichkeit, oder TikTok mit 20\%. Wenn Sie zuerst TikTok benutzen, wechseln Sie danach mit 30\% Wahrscheinlichkeit zu Instagram und bleiben bei TikTok mit 70\% Wahrscheinlichkeit. Wie sieht die Übergangsmatrix aus?
        \end{myexercise}
    }
\end{frame}

\begin{frame}{Die Übergangsmatrix}
    Die Übergangsmatrix fasst alle Übergangswahrscheinlichkeiten zusammen:

    \[
        P =
        \begin{pmatrix}
            0{,}8 & 0{,}2 \\
            0{,}4 & 0{,}6 \\
        \end{pmatrix}
    \]

    \begin{itemize}[<+->]
        \item Zeilen: Ausgangszustände (von welchem Zustand starten wir?)
        \item Spalten: Zielzustände (zu welchem Zustand wechseln wir?)
        \item Jede Zeile summiert sich zu 1 (Wahrscheinlichkeiten!)
    \end{itemize}
\end{frame}

\begin{frame}{Erinnerung: Matrixmultiplikation}
    \only<1-3>{\begin{itemize}[<+->]
            \item Um die Entwicklung einer Markov-Kette zu berechnen, multiplizieren wir einen Zeilenvektor mit einer Matrix.
            \item Beispiel:
        \end{itemize}}
    \[
        \begin{pmatrix}
            \textcolor{red}{a} & \textcolor{blue}{b} \\
        \end{pmatrix}
        \cdot
        \begin{pmatrix}
            \textcolor{teal}{e}   & \textcolor{violet}{f} \\
            \textcolor{orange}{g} & \textcolor{brown}{h}  \\
        \end{pmatrix}
        =
        \begin{pmatrix}
            \textcolor{red}{a} \cdot \textcolor{teal}{e} + \textcolor{blue}{b} \cdot \textcolor{orange}{g} & \textcolor{red}{a} \cdot \textcolor{violet}{f} + \textcolor{blue}{b} \cdot \textcolor{brown}{h} \\
        \end{pmatrix}
    \]
    \only<3->{
        \begin{myexercise}[Matrix-Multiplikation]
            Multiplizieren Sie $\begin{pmatrix} 0{,}3 & 0{,}7 \end{pmatrix}$ mit $\begin{pmatrix} 0{,}8 & 0{,}2 \\ 0{,}4 & 0{,}6 \end{pmatrix}$.
            \only<4->{
                \begin{myanswer}
                    \[
                        \begin{aligned}
                             &
                            \begin{pmatrix}
                                \textcolor{red}{0.3} & \textcolor{blue}{0.7} \\
                            \end{pmatrix}
                            \cdot
                            \begin{pmatrix}
                                \textcolor{teal}{0.8}   & \textcolor{violet}{0.2} \\
                                \textcolor{orange}{0.4} & \textcolor{brown}{0.6}  \\
                            \end{pmatrix}
                            \\&=
                            \begin{pmatrix}
                                \textcolor{red}{0.3} \cdot \textcolor{teal}{0.8} + \textcolor{blue}{0.7} \cdot \textcolor{orange}{0.4} & \textcolor{red}{0.3} \cdot \textcolor{violet}{0.2} + \textcolor{blue}{0.7} \cdot \textcolor{brown}{0.6} \\
                            \end{pmatrix}
                            \\&=
                            \begin{pmatrix}
                                0.24 + 0.28 & 0.06 + 0.42 \\
                            \end{pmatrix}
                            =
                            \begin{pmatrix}
                                0.52 & 0.48 \\
                            \end{pmatrix}
                        \end{aligned}
                    \]
                \end{myanswer}
            }
        \end{myexercise}
    }
\end{frame}

\begin{frame}{Startvektor und Entwicklung}
    \begin{itemize}[<+->]
        \item Startvektor $\pi_0$ gibt den Anfangszustand (=Wahrscheinlichkeit, in einem Zustand zu sein) an
        \item Beispiel: Heute ist es regnerisch
              \[
                  \begin{aligned}
                      \pi_0 =
                       & \begin{pmatrix}
                             0 & 1
                         \end{pmatrix}                                                   \\
                       & \begin{matrix}
                             \;\; \Uparrow & \Uparrow
                         \end{matrix}                                         \\
                       & \begin{matrix}
                             \; \text{\emoji{sun}} & \text{\emoji{umbrella}}
                         \end{matrix}
                  \end{aligned}
              \]
        \item Die Entwicklung nach einem Schritt: $\pi_1 = \pi_0 \cdot P$
    \end{itemize}

    \only<4->{
        \begin{align*}
            \pi_1 = \pi_0 \cdot P & =
            \begin{pmatrix}
                0 & 1 \\
            \end{pmatrix}
            \cdot
            \begin{pmatrix}
                0{,}8 & 0{,}2 \\
                0{,}4 & 0{,}6 \\
            \end{pmatrix}            \\
                                  & =
            \begin{pmatrix}
                0{,}4 & 0{,}6 \\
            \end{pmatrix}
        \end{align*}
    }
\end{frame}

\begin{frame}{Wettervorhersage für übermorgen}
    Wie ist das Wetter in zwei Tagen, wenn es heute regnet?

    \begin{align*}
        \pi_2 = \pi_1 \cdot P & =
        \begin{pmatrix}
            0{,}4 & 0{,}6 \\
        \end{pmatrix}
        \cdot
        \begin{pmatrix}
            0{,}8 & 0{,}2 \\
            0{,}4 & 0{,}6 \\
        \end{pmatrix}            \\
                              & =
        \begin{pmatrix}
            0{,}56 & 0{,}44 \\
        \end{pmatrix}
    \end{align*}

    \only<2->{
        Die Wahrscheinlichkeit, dass es in zwei Tagen sonnig ist, beträgt 56\%.
    }
\end{frame}

\begin{frame}{Auftrag}{Skript}
    \begin{itemize}
        \item \aufgabebox{3.2}
        \item \challengebox{3.3, 3.4}
    \end{itemize}
\end{frame}

\begin{frame}{App-Nutzung: Berechnung}
    Startvektor: $\pi_0 = (1, 0)$ (WhatsApp)

    \begin{align*}
        \pi_1 = \pi_0 \cdot P & =
        \begin{pmatrix}
            1 & 0 \\
        \end{pmatrix}
        \cdot
        \begin{pmatrix}
            0{,}7 & 0{,}3 \\
            0{,}3 & 0{,}7 \\
        \end{pmatrix}            \\
                              & =
        \begin{pmatrix}
            0{,}7 & 0{,}3 \\
        \end{pmatrix}
    \end{align*}

    \begin{align*}
        \pi_2 = \pi_1 \cdot P & =
        \begin{pmatrix}
            0{,}7 & 0{,}3 \\
        \end{pmatrix}
        \cdot
        \begin{pmatrix}
            0{,}7 & 0{,}3 \\
            0{,}3 & 0{,}7 \\
        \end{pmatrix}            \\
                              & =
        \begin{pmatrix}
            0{,}58 & 0{,}42 \\
        \end{pmatrix}
    \end{align*}

    \only<2->{
        Die Wahrscheinlichkeit für Instagram nach zwei Nutzungen beträgt 42\%.
    }
\end{frame}

\begin{frame}{Visualisierung von Markov-Ketten}
    Markov-Ketten können als \textbf{gerichtete Graphen} dargestellt werden:

    \begin{figure}[H]
        \centering
        \begin{tikzpicture}[node distance=3cm, auto, thick]
            \node[state] (S) {Sonnig};
            \node[state] (R) [right of=S] {Regnerisch};

            \path[->]
            (S) edge[bend left] node {0,2} (R)
            (R) edge[bend left] node {0,4} (S)
            (S) edge[loop above] node {0,8} (S)
            (R) edge[loop above] node {0,6} (R);
        \end{tikzpicture}
    \end{figure}

    \begin{itemize}[<+->]
        \item Knoten = Zustände
        \item Kanten = Übergänge mit Wahrscheinlichkeiten
        \item Summe der ausgehenden Kanten = 1
    \end{itemize}
\end{frame}


\begin{frame}{Auftrag: Jupyter Notebooks}
    Skript ``Aus Daten lernen''

    \begin{itemize}
        \item Zuerst jeweils Beispiele genau durchlesen und ausführen
    \end{itemize}
\end{frame}

\begin{frame}{Textvorhersage mit Markov-Ketten}
    \begin{itemize}[<+->]
        \item Markov-Ketten können genutzt werden, um \textbf{Texte zu generieren}.
        \item Die Idee: Jedes Wort (oder Zeichen) ist ein Zustand.
        \item Die Wahrscheinlichkeit für das nächste Wort hängt nur vom aktuellen Wort ab.
    \end{itemize}
    \vspace{0.5cm}
    \begin{block}{Beispiel:}
        \textit{„Ich mag Pizza und ich mag Pasta.“}
    \end{block}
    \vspace{0.3cm}
    \only<3->{
        \begin{itemize}
            \item Übergänge (vereinfacht):\\
                  \texttt{Ich} $\rightarrow$ \texttt{mag} $\rightarrow$ \texttt{Pizza} $\rightarrow$ \texttt{und} $\rightarrow$ \texttt{ich} $\rightarrow$ \texttt{mag} $\rightarrow$ \texttt{Pasta}
        \end{itemize}
    }
\end{frame}

\begin{frame}{Visualisierung: Markov-Kette für Text}
    \centering
    \begin{tikzpicture}[node distance=2.5cm, auto, thick, >=stealth']
        \node[state] (Ich) {\texttt{Ich}};
        \node[state] (mag) [right of=Ich] {\texttt{mag}};
        \node[state] (Pizza) [right of=mag] {\texttt{Pizza}};
        \node[state] (und) [below of=Pizza] {\texttt{und}};
        \node[state] (Pasta) [above of=Pizza] {\texttt{Pasta}};

        \path[->]
        (Ich) edge node {1.0} (mag)
        (mag) edge node[below] {0.5} (Pizza)
        (mag) edge node[below] {0.5} (Pasta)
        (Pizza) edge node {1.0} (und)
        (und) edge node {1.0} (Ich);
    \end{tikzpicture}
    \vspace{0.5cm}

    \only<2->{
        \begin{itemize}
            \item Jeder Pfeil steht für eine mögliche Fortsetzung im Text.
            \item Die Wahrscheinlichkeiten werden aus Beispielsätzen gezählt.
        \end{itemize}
    }
\end{frame}

\begin{frame}[fragile]{Kurze Textgenerierung (Beispiel)}
    \begin{itemize}[<+->]
        \item Start: \texttt{Ich}
        \item Nächstes Wort: \texttt{mag}
        \item Danach: zufällig \texttt{Pizza} oder \texttt{Pasta}
        \item usw.
    \end{itemize}
    \vspace{0.3cm}
    \begin{block}{Mögliche generierte Sätze:}
        \texttt{Ich mag Pizza und ich mag Pasta.}\\
        \texttt{Ich mag Pasta.}
    \end{block}
    \vspace{0.3cm}
    \only<5->{
        Markov-Ketten können einfache, aber oft überraschende Texte erzeugen!
    }
\end{frame}

\begin{frame}
    \begin{itemize}
        \item \aufgabebox{Text-Generierung mit Markov-Ketten}
        \item \challengebox{1.3}
    \end{itemize}
\end{frame}
\begin{frame}{Lernkontrolle: Markov-Ketten für Text}
    \begin{myexercise}

            \texttt{Hallo Welt Info Hallo Info Welt Hallo Hallo Welt Info Info Welt Hallo}

        \begin{enumerate}
            \item Leiten Sie aus dem Text eine Übergangsmatrix für die Wörter \texttt{Hallo}, \texttt{Welt} und \texttt{Info} her.
            \item Was ist der wahrscheinlichste Text mit 3 Wörtern, wenn Sie mit \texttt{Hallo} starten?
        \end{enumerate}
        \textit{Hinweis: Überlegen Sie sich sinnvolle Übergangswahrscheinlichkeiten!}
        
    \end{myexercise}
\end{frame}

\begin{frame}
\begin{enumerate}
            \item \textbf{Übergangszählung:}
            \begin{table}[H]
                                \centering
                                \begin{tabular}{c|ccc}
                                        & \texttt{Hallo} & \texttt{Welt} & \texttt{Info} \\
                                        \hline
                                        \texttt{Hallo} (5x) & 1 & 2 & 1 \\
                                        \texttt{Welt} (4x) & 2 & 0 & 2 \\
                                        \texttt{Info} (4x) & 1 & 2 & 1 \\
                                \end{tabular}
                                \end{table}

                                \textbf{Übergangsmatrix:}

                                Für jede Zeile normieren wir auf die Summe der Übergänge von diesem Wort:

                                \[
                                P = \begin{pmatrix}
                                & \texttt{Hallo} & \texttt{Welt} & \texttt{Info} \\
                                \texttt{Hallo} & \frac{1}{4} & \frac{2}{4} & \frac{1}{4} \\
                                \texttt{Welt} & \frac{2}{4} & 0 & \frac{2}{4} \\
                                \texttt{Info} & \frac{1}{4} & \frac{2}{4} & \frac{1}{4} \\
                                \end{pmatrix}
                                \]

                                \item \textbf{Wahrscheinlichster Text mit 3 Wörtern, Start \texttt{Hallo}:}

                                Von \texttt{Hallo} aus ist der wahrscheinlichste Übergang zu \texttt{Welt} (2 von 4 Übergängen), danach von \texttt{Welt} zu \texttt{Hallo} oder \texttt{Info} (je 2 von 4 Übergängen).

                                \begin{itemize}
                                \item \texttt{Hallo} $\rightarrow$ \texttt{Welt} ($0{,}5$), dann \texttt{Welt} $\rightarrow$ \texttt{Hallo} ($0{,}5$):\\
                                \texttt{Hallo Welt Hallo}
                                \item \texttt{Hallo} $\rightarrow$ \texttt{Welt} ($0{,}5$), dann \texttt{Welt} $\rightarrow$ \texttt{Info} ($0{,}5$):\\
                                \texttt{Hallo Welt Info}
                                \end{itemize}

                                Beide sind gleich wahrscheinlich und am wahrscheinlichsten.

            \textbf{Übergangsmatrix:}

            Für jede Zeile normieren wir auf die Summe der Übergänge von diesem Wort:

            \[
            P = \begin{pmatrix}
            & \texttt{Hallo} & \texttt{Welt} & \texttt{Info} \\
            \texttt{Hallo} & 0 & \frac{3}{6} & \frac{3}{6} \\
            \texttt{Welt} & \frac{2}{5} & 0 & \frac{3}{5} \\
            \texttt{Info} & \frac{3}{5} & \frac{2}{5} & 0 \\
            \end{pmatrix}
            \]

            \item \textbf{Wahrscheinlichster Text mit 3 Wörtern, Start \texttt{Hallo}:}

            Von \texttt{Hallo} aus gibt es gleichwahrscheinliche Übergänge zu \texttt{Welt} und \texttt{Info}. Danach ist der wahrscheinlichste nächste Schritt jeweils der Übergang mit der höchsten Wahrscheinlichkeit:

            \begin{itemize}
            \item \texttt{Hallo} $\rightarrow$ \texttt{Welt} ($0{,}5$), dann \texttt{Welt} $\rightarrow$ \texttt{Info} ($0{,}6$):\\
            \texttt{Hallo Welt Info}
            \item \texttt{Hallo} $\rightarrow$ \texttt{Info} ($0{,}5$), dann \texttt{Info} $\rightarrow$ \texttt{Hallo} ($0{,}6$):\\
            \texttt{Hallo Info Hallo}
            \end{itemize}

            Beide sind gleich wahrscheinlich und am wahrscheinlichsten.
        \end{enumerate}
\end{frame}


\begin{frame}{Grenzen von Markov-Ketten}
    \begin{itemize}[<+->]
        \item \textbf{Markov-Eigenschaft:} Nicht immer realistisch, manchmal hängt die Zukunft von mehr als nur dem aktuellen Zustand ab
        \item \textbf{Feste Übergangswahrscheinlichkeiten:} In der Realität können sich Wahrscheinlichkeiten ändern
        \item \textbf{Komplexe Abhängigkeiten:} Können nicht abgebildet werden
        \item \textbf{Zustandsraum:} Bei vielen Zuständen wird die Berechnung aufwändig
        \item \textbf{Keine Kausalität:} Beschreiben nur Wahrscheinlichkeiten, keine Ursachen
    \end{itemize}
\end{frame}

\begin{frame}[fragile]{Lernziele}
    \begin{itemize}[<+->]
        \item \faCheckSquare{} Ich kann \textbf{Anwendungen} von Markov-Ketten in der Informatik beschreiben.
        \item \faCheckSquare{} Ich kann erklären, was eine \textbf{Markov-Kette} ist und wie die \textbf{Markov-Eigenschaft} funktioniert.
        \item \faCheckSquare{} Ich kann aus einer realen Problemstellung eine \textbf{Übergangsmatrix} erstellen.
        \item \faCheckSquare{} Ich kann mithilfe von Matrizen berechnen, wie sich eine Markov-Kette \textbf{über mehrere Schritte entwickelt}.
              % \item \faSquare{} Ich kann die stationäre Verteilung einer Markov-Kette berechnen und interpretieren.
        \item \faCheckSquare{} Ich kann kurze Texte mit Markov-Ketten \textbf{generieren}.
        \item \faCheckSquare{} Ich kann die \textbf{Grenzen} von Markov-Ketten reflektieren.
    \end{itemize}
\end{frame}

% \begin{frame}
%     \centering
%     \includegraphics[width=0.7\textwidth]{"GrundlagenInfo/10_AusDatenLernen/Figures/markov-meme-chatGPT-2"}
% \end{frame}

% \begin{frame}
%     \centering
%     \includegraphics[width=0.7\textwidth]{"GrundlagenInfo/10_AusDatenLernen/Figures/markov-meme-chatGPT-3"}
% \end{frame}

% \begin{frame}
%     \centering
%     \includegraphics[width=0.7\textwidth]{"GrundlagenInfo/10_AusDatenLernen/Figures/markov-meme-chatGPT-1"}
% \end{frame}

\appendix
\begin{frame}[plain,noframenumbering]
    \centering
    \vspace{2cm}
    {\Huge\textbf{Anhang}}
\end{frame}

\begin{frame}{Eigenschaften von Markov-Ketten}
    \begin{columns}
        \column{0.5\textwidth}
        \begin{itemize}[<+->]
            \item \textbf{Irreduzibilität:} Von jedem Zustand kann man jeden anderen erreichen
            \item \textbf{Aperiodizität:} Rückkehr zu einem Zustand nicht nur in festen Zeitabständen
            \item \textbf{Absorbierende Zustände:} Zustände, die nicht mehr verlassen werden
        \end{itemize}

        \column{0.5\textwidth}
        \begin{itemize}[<+->]
            \item \textbf{Rekurrenz:} Wahrscheinlichkeit, wieder in einen Zustand zurückzukehren
            \item \textbf{Transienz:} Zustände, die irgendwann nicht mehr erreicht werden
            \item \textbf{Ergodizität:} Positive Rekurrenz und Aperiodizität
        \end{itemize}
    \end{columns}
\end{frame}

\begin{frame}{Absorbierende Zustände}
    \begin{columns}
        \column{0.6\textwidth}
        Beispiel:
        \[
            P = \begin{pmatrix}
                1     & 0     \\
                0{,}3 & 0{,}7 \\
            \end{pmatrix}
        \]

        \column{0.4\textwidth}
        \begin{figure}[H]
            \centering
            \begin{tikzpicture}[node distance=2cm, auto, thick, scale=0.7, transform shape]
                \node[state] (1) {1};
                \node[state] (2) [below of=1] {2};

                \path[->]
                (1) edge[loop right] node {1,0} (1)
                (2) edge[bend left] node {0,3} (1)
                (2) edge[loop right] node {0,7} (2);
            \end{tikzpicture}
        \end{figure}
    \end{columns}

    \only<2->{
        \begin{itemize}
            \item Zustand 1 ist \textbf{absorbierend}: $P_{1,1} = 1$
            \item Wenn die Kette in Zustand 1 ist, bleibt sie dort für immer
            \item Zustand 2 ist \textbf{transient}: Irgendwann wird er verlassen und nie wieder erreicht
        \end{itemize}
    }
\end{frame}

\begin{frame}{Stationäre Verteilung}
    \begin{mydefinition}[Stationäre Verteilung]
    Eine Wahrscheinlichkeitsverteilung $\pi = (\pi_1, \pi_2, \ldots)$ heisst \textbf{stationäre Verteilung} einer Markov-Kette mit Übergangsmatrix $P$, wenn gilt:
    \[
        \pi = \pi \cdot P
    \]
    \end{mydefinition}

    \begin{itemize}[<+->]
        \item Langfristige, stabile Verteilung der Wahrscheinlichkeiten
        \item Bleibt nach Multiplikation mit $P$ unverändert
        \item Für endliche, irreduzible, aperiodische Markov-Ketten immer eindeutig
        \item Unabhängig vom Anfangszustand erreicht
    \end{itemize}
\end{frame}

\begin{frame}{Berechnung der stationären Verteilung}
    Wetter-Beispiel: $P = \begin{pmatrix} 0{,}8 & 0{,}2 \\ 0{,}4 & 0{,}6 \end{pmatrix}$

    \begin{enumerate}[<+->]
        \item $\pi = \pi \cdot P$ aufstellen: $(\pi_1, \pi_2) = (\pi_1, \pi_2) \cdot \begin{pmatrix} 0{,}8 & 0{,}2 \\ 0{,}4 & 0{,}6 \end{pmatrix}$
        \item Gleichungen: $\pi_1 = 0{,}8\pi_1 + 0{,}4\pi_2$ und $\pi_2 = 0{,}2\pi_1 + 0{,}6\pi_2$
        \item Umformen: $0{,}2\pi_1 = 0{,}4\pi_2$ → $\pi_1 = 2\pi_2$
        \item Normierung: $\pi_1 + \pi_2 = 1$ → $2\pi_2 + \pi_2 = 1$ → $3\pi_2 = 1$ → $\pi_2 = \frac{1}{3}$
        \item Ergebnis: $\pi_1 = \frac{2}{3}$ und $\pi_2 = \frac{1}{3}$
    \end{enumerate}

    \only<6->{
        Langfristig ist es an 2/3 der Tage sonnig und an 1/3 der Tage regnerisch.
    }
\end{frame}

\begin{frame}{App-Nutzung: Stationäre Verteilung}
    Für die App-Nutzung: $P = \begin{pmatrix} 0{,}7 & 0{,}3 \\ 0{,}3 & 0{,}7 \end{pmatrix}$

    \begin{enumerate}[<+->]
        \item Aus $\pi_1 = 0{,}7\pi_1 + 0{,}3\pi_2$ folgt $0{,}3\pi_1 = 0{,}3\pi_2$ → $\pi_1 = \pi_2$
        \item Mit $\pi_1 + \pi_2 = 1$ folgt $\pi_1 = \pi_2 = 0{,}5$
    \end{enumerate}

    \only<3->{
        Langfristig werden beide Apps gleich häufig genutzt.
    }
\end{frame}